\documentclass[12pt,conference,onecolumn]{IEEEtran}

\usepackage[hidelinks]{hyperref}

\title{Closed-loop music therapy for pain management using heart rate variability ({HRV}) and patient-specific musical preference}
\author{%
\IEEEauthorblockN{Sarah Afzal}\IEEEauthorblockA{Science \& Engineering\\Manalapan High School\\Englishtown, NJ\\\href{mailto:427safzal@frhsd.com}{427safzal@frhsd.com}}\and
\IEEEauthorblockN{Saanvi Dakwale}\IEEEauthorblockA{Science \& Engineering\\Manalapan High School\\Englishtown, NJ\\\href{mailto:427sdakwale@frhsd.com}{427sdakwale@frhsd.com}}\and
\IEEEauthorblockN{K August}\IEEEauthorblockA{\href{mailto:kit.august@gmail.com}{kit.august@gmail.com}}
}

\date{January 28, 2027}

\newcommand{\keywords}{IEEE, internship, biofeedback, heart-rate variability, music, breathing, biomedical signal processing, pain management}

\usepackage{hyperref}
\makeatletter
\AtBeginDocument{
\hypersetup{%
pdftitle={\@title},
pdfauthor={Sarah Afzal, Saanvi Dakwale, and K M August},
pdfkeywords={\keywords}}}
\makeatother

\begin{document}
\maketitle 

\begin{abstract}
Pain produces measurable changes in the autonomic nervous system, reflected in heart rate variability (HRV) and related cardiovascular signals, shifting the body away from a calm, parasympathetic-dominant state. This project investigates whether music combined with pace-feeding biofeedback to maintain HRV using a low-cost Photoplethysmograph (PPG) sensor before, during, and after a mild, cold-water pain stimulus on 8 participants. On day 1, participants will have a song randomly selected to them and on day two, they will pick the song from the same four options, where they will manually match their breathing to the tempo of the music. The PPG signals will then undergo signal processing to calculate and quantify HRV through measurements of RMSSD and SDNN, as well as identify ectopic beats. It is hypothesized that music combined with paced-breathing biofeedback will produce greater increases in HRV and faster recovery than the no-music condition. This project aims to execute a design study to develop a closed-loop prototype integrating low-cost physiological sensing, signal processing, quantitative music analysis, and real-time biofeedback to investigate whether musical preference and physiological matching can contribute to measurable autonomic regulation during pain.
\end{abstract}

\begin{IEEEkeywords}
\keywords
\end{IEEEkeywords}

\end{document}

%Sources:
%Chen, S., Yuan, Q., Wang, C., Ye, J., & Yang, L. (2025). The effect of music therapy for patients with chronic pain: systematic review and meta-analysis. BMC psychology, 13(1), 455. https://doi.org/10.1186/s40359-025-02643-x
%de Witte, M., Pinho, A. da S., Stams, G. J., Moonen, X., Bos, A. E. R., & van Hooren, S. (2022). Music therapy for stress reduction: a systematic review and meta-analysis. Health Psychology Review, 16(1), 134–159. https://doi.org/10.1080/17437199.2020.1846580
%Getie, A., Bimerew, M. & Kitaw, T.A. The effect of music therapy on pain management: a systematic review and meta-analysis of randomized controlled trials. Discov Psychol 6, 54 (2026). https://doi.org/10.1007/s44202-026-00600-2
%Raglio, A., Attardo, L., Gontero, G., Rollino, S., Groppo, E., & Granieri, E. (2015). Effects of music and music therapy on mood in neurological patients. World journal of psychiatry, 5(1), 68–78. https://doi.org/10.5498/wjp.v5.i1.68
